\documentclass[11pt,a4paper]{article}
\usepackage{amsmath, amssymb, amsfonts}
\usepackage{geometry}
\usepackage{hyperref}
\geometry{margin=1in}
\usepackage{parskip}
\usepackage{xcolor}

\newcommand{\R}{\mathbb{R}}

\begin{document}
	
	\title{\textbf{EngageNet Draft}}
	\author{Lado Turmanidze}
	\date{May 4, 2026}
	\maketitle
	
	\section{Architecture}
	\subsection{Input}
	Let $x=\{x_{1},\dots,x_{M}\}$ with $x_{i}\in\R^{C_{i}\times L}$ for all $i\in\{1,\dots,M\}$, where $C_{i}$ is the number of acquisition channels for the $i$-th modality and $L$ represents the number of time steps of the input signal.
	
	\subsection{Output}
	In \cite{ziyu2025}, $y$ is a discrete value corresponding to a category, which is then converted into a one-hot encoding $p(y|x)\in\R^{N}$ (where $N$ is the number of categories). In our case, since we are performing regression, $y\in[0,1] \implies p(y|x)\in[0,1]^{N}$, so our model outputs the predicted category probability $p(\hat{y}|x)\in[0,1]^{N}$.
	
	\subsection{Inner Workings}
	We need to capture contextual information from physiological time series simultaneously, and a unidirectional approach cannot fully capture the complex temporal patterns. The multimodal BiMamba network consists of two modules: the intra-modal BiMamba module and the inter-modal BiMamba module. The former is designed to model each modality independently, while the latter effectively fuses multimodal features.
	
	\subsubsection{Intra-modal BiMamba Module}
	First, an initial encoder $h_{i} = \text{InitEncoder}_{i}(x_{i}) \in \R^{L^{\prime}\times C_{i}^{\prime}}$ for all $i\in\{1,\dots,M\}$ is used to extract shallow features from a specific modality, where $C_{i}^{\prime}$ is the number of output channels and $L^{\prime}$ is the feature length of the output signal. 
	
	The transposition of dimensions from $C_i \times L$ to $L^{\prime} \times C_i^{\prime}$ physically aligns the data structure with the sequence-first requirement of autoregressive models. By enforcing the sequence length $L^{\prime}$ as the primary operational dimension, the subsequent state-space and 1D convolutional operators can process the temporal dimension across the channel feature space.
	
	The BiMamba module operates as follows:
	\begin{enumerate}
		\item The output of the gating mechanism is $g_{i}=\sigma(W_{i}^{g}h_{i}+b_{i}^{g})$, where $h_{i}$ is the hidden state, $W_{i}^{g}$ is the weight matrix, $b_{i}^{g}$ is the bias, and $\displaystyle \sigma(x)=x\cdot\frac{1}{1+e^{-x}}$ (SiLU). Intuitively, $g_{i}$ highlights emotion-relevant information while suppressing noise and redundancy.
		
		\item BiMamba models temporal dependencies from both forward and backward perspectives using state-space modeling, thereby capturing richer contextual dynamics:
		$$h_{i}^{\rightarrow}=g_{i}\otimes \text{SSM}_{\rightarrow}(\sigma(\text{Conv1D}_{\rightarrow}(W_{i}^{h}h_{i}+b_{i}^{h})))\in\R^{L^{\prime}\times C_{i}^{\prime}}$$
		$$h_{i}^{\leftarrow}=g_{i}\otimes \text{SSM}_{\leftarrow}(\sigma(\text{Conv1D}_{\leftarrow}(\text{Rev}_{t}(W_{i}^{h}h_{i}+b_{i}^{h}))))\in\R^{L^{\prime}\times C_{i}^{\prime}}$$
		where $h_{i}^{\rightarrow}$ and $h_{i}^{\leftarrow}$ are the feature representations of the $i$-th modality in forward and backward modeling. $\text{Rev}_{t}(\cdot)$ denotes flipping along the time dimension $L^{\prime}$.
		
		The operator $\otimes$ denotes the Hadamard product (element-wise multiplication), which explains why the dimensions remain strictly unchanged, unlike a direct sum ($\oplus$) or tensor product. $W_{i}^{h}$ represents the weight matrix and $b_{i}^{h}$ represents the bias. 
		
		$\text{SSM}(\cdot)$ represents the Selective State Space Model. Formally, it maps a 1-dimensional sequence through a continuous linear time-invariant (LTI) system defined by the state equation $x'(t) = A x(t) + B u(t)$ and output $y(t) = C x(t)$. To operate on discrete sequence steps, the continuous transition operator $A$ and input operator $B$ are discretized via a timescale parameter $\Delta$, yielding the discrete operators $\bar{A} = \exp(\Delta A)$ and $\bar{B} = (\Delta A)^{-1}(\bar{A} - I)\Delta B$. 
		
		Unlike traditional LTI systems, Mamba introduces \textit{selection} by making the discretized parameters $(\bar{B}, C, \Delta)$ functions of the input $u_t$. Specifically:
		$$ \Delta_t = \text{softplus}(s_{\Delta} + \text{Linear}(u_t)) $$
		$$ B_t = \text{Linear}(u_t), \quad C_t = \text{Linear}(u_t) $$
		where $ \text{softplus}(x) = \ln(1 + e^x) \text{ and } s_{\Delta} \in \mathbb{R}^{D} $ is a learned, input-independent bias term that defines the baseline timescale of the system. The inclusion of $s_{\Delta}$ within the \textit{softplus} function ensures that the resulting $\Delta_t$ remains strictly positive, which is required for the stability of the discretization. This selection mechanism allows the discrete transition $x_t = \bar{A}_t x_{t-1} + \bar{B}_t u_t$ to dynamically adjust its temporal resolution - effectively allowing the model to "reset" its state or focus on transient patterns based on the specific content of the sequence $u_t$.
		
		\item BiMamba applies a linear projection and a  \hypertarget{residual}{\textcolor{red}{residual connection}} to integrate bidirectional features and stabilize training:
		$$u_{i}=h_{i}+\left(W_{i}^{o}\left(\frac{h_{i}^{\rightarrow}+\text{Rev}_{t}(h_{i}^{\leftarrow})}{2}\right)+b_{i}^{o}\right)\in\R^{L^{\prime}\times C_{i}^{\prime}}$$
		which is the output of the intra-modal BiMamba module for the $i$-th modality. $W_{i}^{o}$ is the weight matrix, and $b_{i}^{o}$ is the bias.
	\end{enumerate}
	
	\subsubsection{Inter-modal BiMamba Module}
	Relying solely on intra-modal modeling neglects inter-modal correlations that are crucial for comprehensive multimodal understanding. Let the concatenation along the channel dimension be denoted by the operator $\circ$. We define the concatenated intermediate matrix $U$:
	$$U = u_1 \circ u_2 \circ \dots \circ u_M$$
	where $U \in \R^{L^{\prime} \times \sum_{i=1}^{M} C_i^{\prime}}$. To enable state-space modeling across the modalities rather than the temporal space, we transpose this matrix:
	$$m = U^{\top} = \text{Transpose}(u_1 \circ u_2 \circ \dots \circ u_M) \in \R^{(\sum C_i^{\prime}) \times L^{\prime}}$$
	Thus, each time step acts structurally as a "channel," while the concatenated modality dimension serves as the sequence length. We then input $m$ into a BiMamba structure:
	$$ H = \text{BiMamba}(m) \in \R^{\sum_{i=1}^{M} C_i^{\prime} \times L^{\prime}} $$ By performing bidirectional state modeling along the fused $\sum_{i=1}^{M} C_i^{\prime}$ axis, the sequence transition operators effectively treat preceding modalities as the context basis for subsequent modalities.
	
	The function $H = \text{BiMamba}(m)$ represents the application of the entire bidirectional Mamba block (All 1--3 steps of the intra-modal module above) to the fused matrix $m$. In this inter-modal stage, the "sequence length" is now the total number of concatenated channels $\sum C_i^{\prime}$, and the "features" at each step are the temporal points $L^{\prime}$. This allows the model to build a context where the state of one modality can directly influence the hidden state of another.
	
	In the forward process, the hidden states built from modalities $\{1, \dots, i- 1 \}$, while in the backward process the hidden states derived from the modalities $\{M, M - 1, \dots, i + 1\}$ provide additional temporal information for the $i$-th modality.
	
	\section{Possible Improvement Ideas}
	\begin{enumerate}
		\item \hyperlink{residual}{Residual connection}: we can replace the standard residual connection with DeepSeek AI's Manifold-Constrained Hyper-Connections (mHC) \cite{xie2026mhc}. 
		
		Let the BiMamba transformation block be denoted as the layer function $$ \mathcal{F}(h_i) = W_{i}^{o}\left(\frac{h_{i}^{\rightarrow}+\text{Rev}_{t}(h_{i}^{\leftarrow})}{2}\right)+b_{i}^{o} $$ In the original architecture, the (residual) update is $u_i = h_i + \mathcal{F}(h_i) $. 
		
		To apply mHC, we expand the residual stream width by a factor of $n$ to form a multi-stream representation $\mathbf{h}_i \in \R^{n \times L^{\prime} \times C_i^{\prime}}$. The update rule becomes:
		$$u_i = \mathcal{H}_i^{res} \mathbf{h}_i + (\mathcal{H}_i^{post})^{\top} \mathcal{F}(\mathcal{H}_i^{pre} \mathbf{h}_i)$$
		where $\mathcal{H}_i^{pre} \in \R^{1 \times n}$ aggregates the multi-stream into a single input for $\mathcal{F}$, $\mathcal{H}_i^{post} \in \R^{1 \times n}$ maps the output back to the stream, and $\mathcal{H}_i^{res} \in \R^{n \times n}$ mixes features within the stream. 
		
		The structural mappings $\mathcal{H}_i^{pre}$ and $\mathcal{H}_i^{post}$ are composed of dynamic, input-dependent linear projections and static, learnable biases. To prevent signal cancellation, they are mapped through a \textit{sigmoid} activation $\sigma$:
		
		$$\mathcal{H}_i^{pre} = \sigma(\alpha_{i}^{pre} (\vec{\mathbf{h}}_i \varphi_{i}^{pre}) + b_{i}^{pre})$$
		
		$$\mathcal{H}_i^{post} = 2\sigma(\alpha_{i}^{post} (\vec{\mathbf{h}}_i \varphi_{i}^{post}) + b_{i}^{post})$$
		
		where $\vec{\mathbf{h}}_i$ is the flattened input tensor, $\varphi$ represents the linear projection weights, and $\alpha$ acts as a learnable gating factor.
		
		To maintain signal conservation, $\mathcal{H}_i^{res}$ is projected onto the Birkhoff polytope using the Sinkhorn-Knopp algorithm. The Birkhoff polytope is the convex hull of all permutation matrices, consisting entirely of doubly stochastic matrices (square matrices with non-negative real entries where each row and column sums exactly to 1). 
		
		The projection is achieved iteratively:
		$$\mathcal{H}_i^{res} = \text{Sinkhorn-Knopp}(\tilde{\mathcal{H}}_i^{res})$$
		Given an initial positive matrix $M^{(0)} = \exp(\tilde{\mathcal{H}}_i^{res})$, the Sinkhorn-Knopp algorithm alternately normalizes the rows and columns until convergence. Constraining $\mathcal{H}_i^{res}$ to the Birkhoff polytope ensures its spectral norm is strictly bounded by 1 ($||\mathcal{H}_i^{res}||_2 \le 1$). This property guarantees stability across deep layers, strictly regulating the vector norm and preventing the signal from exploding or vanishing within the Hilbert space of our feature representations.
	\end{enumerate}
	
	\section{Engagement Regression Adaptation}
	
	\subsection{Differential Entropy and the Failure of Gaussian Assumptions}
	To adapt the entropy-based selection mechanism of BiMamba-TTA for continuous engagement regression, we transition from Shannon entropy to \textbf{differential entropy}. For a continuous random variable $y$ with probability density $p(y)$, the differential entropy is defined as:
	$$ H(y) = -\int_{\mathcal{Y}} p(y) \ln p(y) \, dy $$
	While a Gaussian assumption $p(y|x) = \mathcal{N}(\mu, \sigma^2)$ simplifies this to $H(y) = \frac{1}{2}\ln(2\pi e \sigma^2)$, it is fundamentally inappropriate for engagement modeling. Engagement scores are strictly bounded within the interval $[0, 1]$. A Gaussian distribution assigns non-zero probability mass to values in $(-\infty, 0)$ and $(1, \infty)$, leading to "probability leakage" and poor calibration near the boundaries.
	
	\subsection{Beta Regression Framework}
	To respect the support $\mathcal{Y} \in [0, 1]$, we model the predictive distribution using a \textbf{Beta distribution}:
	$$ p(y|\alpha, \beta) = \frac{y^{\alpha-1}(1-y)^{\beta-1}}{B(\alpha, \beta)}, \quad B(\alpha, \beta) = \frac{\Gamma(\alpha)\Gamma(\beta)}{\Gamma(\alpha+\beta)} $$
	The network outputs these parameters via a softplus activation: $\alpha = \text{softplus}(z_\alpha) + 1$ and $\beta = \text{softplus}(z_\beta) + 1$. The $+1$ offset ensures the distribution remains unimodal ($\alpha, \beta > 1$). 
	
	The predictive mean $\hat{\mu}$ and aleatoric uncertainty $\text{Var}[y]$ are:
	$$ \hat{\mu} = \frac{\alpha}{\alpha + \beta}, \quad \text{Var}[y] = \frac{\alpha\beta}{(\alpha+\beta)^2(\alpha+\beta+1)} $$
	The Beta variance is naturally boundary-aware, satisfying: 
	$$ \lim_{\hat{\mu} \to 0} \text{Var}[y] = 0 \quad \text{and} \quad \lim_{\hat{\mu} \to 1} \text{Var}[y] = 0 $$
	This provides a principled measure of confidence that scales with the proximity to the interval limits.
	
	\subsection{Test-Time Adaptation (TTA) and Adaptive Thresholds}
	
	Following the differential-entropy principle of the original BiMamba-TTA paper \cite{ziyu2025}, we replace Shannon entropy with a regression-compatible uncertainty proxy.
	For the Beta predictive distribution $p(y \mid x) = \mathrm{Beta}(\hat{\alpha}(x),\hat{\beta}(x))$ that respects the support $y \in [0,1]$, we define
	$$
	U(x) := \mathrm{Var}[y] = \frac{\hat{\alpha}(x)\hat{\beta}(x)}{(\hat{\alpha}(x)+\hat{\beta}(x))^2(\hat{\alpha}(x)+\hat{\beta}(x)+1)}.
	$$
	
	\subsubsection{Why Predictive Variance is a Strictly Monotonic Proxy for Differential Entropy}
	
	This choice is justified because variance is a \emph{strictly monotonic} function of the differential entropy $H(p)$ of the Beta distribution.  
	To see this explicitly, re-parameterise the Beta distribution by its mean and concentration:
	$$
	\mu = \frac{\hat{\alpha}}{\hat{\alpha}+\hat{\beta}}, \qquad \kappa = \hat{\alpha} + \hat{\beta}.
	$$
	Substituting these parameters yields the simple expression
	$$
	U(x) = \frac{\mu(1-\mu)}{\kappa + 1}.
	$$
	For any fixed mean $\mu \in (0,1)$, the partial derivative
	$$
	\frac{\partial U}{\partial \kappa} = -\frac{\mu(1-\mu)}{(\kappa+1)^2} < 0
	$$
	shows that variance is strictly decreasing in $\kappa$: larger concentration $\kappa$ produces a tighter distribution and therefore smaller variance.
	
	The differential entropy of the Beta distribution has the closed-form expression
	$$
	H(p) = \ln B(\hat{\alpha},\hat{\beta}) - (\hat{\alpha}-1)\psi(\hat{\alpha}) - (\hat{\beta}-1)\psi(\hat{\beta}) + (\hat{\alpha}+\hat{\beta}-2)\psi(\hat{\alpha}+\hat{\beta}),
	$$
	where $B(\cdot,\cdot)$ is the Beta function and $\psi(\cdot)$ is the digamma function.  
	Substituting the same re-parameterisation $\hat{\alpha}=\mu\kappa$, $\hat{\beta}=(1-\mu)\kappa$ and taking the derivative with respect to $\kappa$ likewise yields
	$$
	\frac{\partial H}{\partial \kappa} < 0
	$$
	for any fixed $\mu \in (0,1)$. Thus both $U(x)$ and $H(p)$ are strictly decreasing functions of $\kappa$: larger concentration (tighter distribution) produces both smaller variance \emph{and} smaller differential entropy.  
	Consequently, the ranking of samples by uncertainty is identical whether one uses the full (expensive) differential entropy expression or simply the predictive variance.
	
	Variance is numerically stable because its formula consists only of basic arithmetic operations (no logarithms, no digamma or Gamma functions, and no risk of underflow when $\hat{\alpha}$ or $\hat{\beta}$ become large).  
	It is boundary-aware: $\lim_{\mu\to 0^+}U=0$ and $\lim_{\mu\to 1^-}U=0$, which matches the intuition that a prediction that is certain to be near the extremes of engagement (fully disengaged or fully engaged) should carry low uncertainty.  
	Finally, variance is already required for the Beta negative log-likelihood (Beta NLL) loss used during training:
	$$
	\mathcal{L}_{\text{NLL}}(x, y) = -[(\hat{\alpha}-1)\ln y + (\hat{\beta}-1)\ln(1-y) - \ln B(\hat{\alpha},\hat{\beta})].
	$$
	Because the variance appears directly in the predictive distribution (and therefore in every forward pass), it incurs no extra computation at inference time.  
	All these properties make variance the natural, cheap uncertainty proxy for regression.
	
	A sample is selected for the TTA backward pass if the multimodal prediction is confident (low uncertainty) but the unimodal branches are uncertain (high uncertainty), i.e.,
	$$
	S(x) = \mathbb{I}\Bigl\{ U_{\text{multi}}(x) \leq \gamma_m(t) \;\text{AND}\; \sum_{i=1}^M w_i U_i(x) \geq \gamma_u(t) \Bigr\}.
	$$
	
	To handle non-stationary streaming data (i.e., test-time engagement signals that arrive continuously frame-by-frame and whose statistical distribution drifts over time due to fatigue, cultural differences, sensor drift, or changing speaker roles), the thresholds $\gamma(t)$ are updated adaptively using a moving window of the most recent $K$ samples:
	$$
	\gamma(t) = \mathrm{Percentile}\bigl(\mathcal{U}_t,\, p\bigr), \quad
	\mathcal{U}_t = \bigl\{U(\tau) \mid \tau \in [t-K,t]\bigr\},
	$$
	where $p$ is the percentile rank (e.g., $p=0.8$) that defines the "confidence" cutoff relative to the recent noise floor of the data stream.  
	Here "non-stationary streaming data" means the incoming test distribution is not fixed but changes gradually; the moving window $\mathcal{U}_t$ tracks this drift in real time.  
	The "percentile rank $p=0.8$" selects the value below which $80\,\%$ of the recent uncertainties lie, so the threshold automatically adapts to the current level of noise.  
	The term "confidence cutoff" means that samples with $U_{\text{multi}}(x)$ below this adaptive value are considered sufficiently reliable for adaptation.  
	The "recent noise floor" refers to the lowest typical uncertainty observed in the immediate past window - using the percentile of this floor makes the threshold robust to gradual changes in signal quality.
	
	\subsection{Uncertainty Estimation Methods}
	
	Uncertainty is formally quantified as the variance of the predictive distribution. We utilize three complementary methods that together capture both aleatoric and epistemic uncertainty. All three are expressed in terms of the Beta predictive distribution $p(y \mid x) = \mathrm{Beta}(\hat{\alpha}(x),\hat{\beta}(x))$.
	
	\begin{enumerate}
		\item \textbf{Parametric Beta (Aleatoric Uncertainty):}  
		This is the direct predictive variance obtained from the dual-head Beta output:
		$$
		U(x) = \mathrm{Var}[y] = \frac{\hat{\alpha}(x)\hat{\beta}(x)}{(\hat{\alpha}(x)+\hat{\beta}(x))^2(\hat{\alpha}(x)+\hat{\beta}(x)+1)}.
		$$
		Because the variance is computed analytically from the learned shape parameters, it is the fastest and default choice for the TTA filter.
		
		\item \textbf{Monte Carlo Dropout (Epistemic Uncertainty):}  
		We keep dropout active at inference and draw $T$ stochastic forward passes (typically $T=10$--$20$). For each pass $t$ we obtain a Beta distribution with parameters $(\hat{\alpha}^{(t)},\hat{\beta}^{(t)})$ and corresponding mean $\displaystyle \hat{\mu}^{(t)} = \frac{\hat{\alpha}^{(t)}}{\hat{\alpha}^{(t)}+\hat{\beta}^{(t)}}$.  
		The epistemic uncertainty is the variance of these means:
		$$
		U_{\text{epi}}(x) = \frac{1}{T}\sum_{t=1}^T \bigl(\hat{\mu}^{(t)}(x) - \bar{\mu}(x)\bigr)^2,
		$$
		where $\bar{\mu}(x) = \frac{1}{T}\sum_{t=1}^T \hat{\mu}^{(t)}(x)$. This term quantifies the model's "lack of knowledge" (how sensitive the prediction is to random weight masks) and is especially useful under distribution shift (NoXi $\to$ NoXi+J).
		
		\item \textbf{Total Uncertainty (Hybrid):}  
		Using the law of total variance, the full predictive uncertainty decomposes as
		$$
		U_{\text{total}}(x) = \underbrace{\frac{1}{T}\sum_{t=1}^T \mathrm{Var}\bigl[\mathrm{Beta}(\hat{\alpha}^{(t)},\hat{\beta}^{(t)})\bigr]}_{\text{aleatoric (data)}} 
		+ \underbrace{\frac{1}{T}\sum_{t=1}^T \bigl(\hat{\mu}^{(t)}(x) - \bar{\mu}(x)\bigr)^2}_{\text{epistemic (model)}}.
		$$
		The decomposition is mainly used as a diagnostic to monitor whether the shift is mostly epistemic.
	\end{enumerate}
	
	Note that the symbol $U(x)$ in the TTA filter always refers to one of the three quantities above (most often the parametric Beta variance for speed or the hybrid total variance when epistemic information is needed). The law of total variance guarantees that $U_{\text{total}}$ is the correct variance of the mixture distribution induced by MC dropout.
	
	\subsection{Test-Time Loss Derivations}
	
	The TTA objective enforces consistency between the unimodal student heads and the multimodal teacher while still supervising the students toward the multimodal prediction. We define the \textbf{Mutual Information Sharing (MIS) loss} as the KL divergence that pushes each unimodal Beta distribution toward the multimodal one:
	$$
	\mathcal{L}_{\text{mis}} = \sum_{i=1}^M D_{KL} \bigl( \mathrm{Beta}(\alpha_i, \beta_i) \parallel \mathrm{Beta}(\alpha_{\text{multi}}, \beta_{\text{multi}}) \bigr),
	$$
	where the closed-form Beta KL divergence, with $ B_i := B(\alpha_i,\beta_i) $, is:
	$$
	D_{KL}(B_1 \parallel B_2) = \ln \frac{B(\alpha_2,\beta_2)}{B(\alpha_1,\beta_1)} + (\alpha_1 - \alpha_2)\psi(\alpha_1) + (\beta_1 - \beta_2)\psi(\beta_1) + (\alpha_2 - \alpha_1 + \beta_2 - \beta_1)\psi(\alpha_1 + \beta_1)
	$$
	
	The final TTA loss combines the MIS consistency term with a supervision term that treats the multimodal mean as a pseudo-target for each unimodal head:
	$$
	\mathcal{L}_{TTA} = \mathcal{L}_{\text{mis}} + \lambda \sum_{i=1}^M \text{NLL}_{\text{Beta}}\bigl(\hat{\mu}_{\text{multi}},\, \alpha_i, \beta_i\bigr).
	$$
	Here $\lambda > 0$ is a hyperparameter that controls the relative importance of the two terms: larger $\lambda$ emphasises direct supervision from the multimodal teacher, while smaller $\lambda$ puts more weight on cross-modal consistency. The term $\text{NLL}_{\text{Beta}}(\hat{\mu}_{\text{multi}}, \alpha_i, \beta_i)$ is the Beta negative log-likelihood evaluated at the multimodal mean $\hat{\mu}_{\text{multi}}$ (treated as a pseudo-label) with the unimodal shape parameters.
	
	\subsection{Surgical Fine-Tuning and Variance Starvation}
	
	We only update specific surgical layers to prevent catastrophic forgetting:
	\begin{itemize}
		\item \textbf{BatchNorm layers:} These encode per-channel activation statistics, allowing for immediate adaptation to domain shifts.
		\item \textbf{First convolution in InitEncoder:} Re-calibrates low-level feature extraction (e.g., sensor noise) before temporal/modal mixing.
		\item \textbf{First FC layer in Inter-modal BiMamba:} Recalibrates cross-modal fusion (e.g., adjusting the weight of audio vs. visual cues for different cultures).
	\end{itemize}
	
	This surgical approach mitigates the \textbf{Variance Starvation Problem (VSP)}. In VSP, a model minimises the Beta NLL by artificially inflating $\mathrm{Var}[y]$ (i.e., reducing the concentration $\kappa = \hat{\alpha}+\hat{\beta}$) on hard samples; the gradient to the mean head then vanishes because the loss term becomes insensitive to prediction error. By freezing the majority of the backbone and updating only the listed surgical layers, we force the model to improve genuine feature alignment rather than using uncertainty as a "slack variable" to bypass difficult data.
	
	\section{Permutation-Equivariant Inter-modal Fusion}
	
	\subsection{The Order-Sensitivity Problem}
	
	Recall from Section 1.2.2 that the inter-modal BiMamba module processes the transposed concatenation
	$$m = \text{Transpose}(u_1 \circ u_2 \circ \dots \circ u_M) \in \R^{(\sum_{i=1}^{M} C_i') \times L'},$$
	treating the fused channel axis as a sequence of length $\sum C_i'$, with each modality's channels occupying a contiguous block. The ordering of these blocks - audio before video, or video before pose - is chosen arbitrarily by the practitioner; no prior knowledge assigns a canonical rank to distinct measurement modalities. We show that this arbitrariness introduces a structural bias, because selective state-space models are inherently order-sensitive.
	
	Let $\xi_k \in \R^{L'}$ denote the SSM hidden state at sequence position $k$, and let $s_k \in \R^{L'}$ be the input at that position (corresponding to the channel block of the $k$-th modality in $m$). The Mamba selective SSM \cite{gu2023mamba} defines the recurrence
	$$\xi_k = \bar{A}_k(s_k)\,\xi_{k-1} + \bar{B}_k(s_k)\,s_k, \qquad y_k = C_k(s_k)\,\xi_k,$$
	where $\xi_0 = \mathbf{0}$ and $\bar{A}_k,\bar{B}_k,C_k$ are input-dependent operators instantiated via the selection mechanism.
	
	\textbf{Proposition.} For any SSM with non-trivial state transition (i.e.\ $\bar{A} \neq I$), the output sequence is not invariant under permutations of the input.
	
	\textit{Proof.} It suffices to consider $M = 2$ and a linear, non-selective SSM with constant operators $\bar{A}, \bar{B}, C$. Under the ordering $(u_1, u_2)$:
	$$\xi_1 = \bar{B}\,u_1, \qquad y_2 = C\!\left(\bar{A}\bar{B}\,u_1 + \bar{B}\,u_2\right).$$
	Under the swapped ordering $(u_2, u_1)$:
	$$\xi_1' = \bar{B}\,u_2, \qquad y_2' = C\!\left(\bar{A}\bar{B}\,u_2 + \bar{B}\,u_1\right).$$
	The difference is
	$$y_2 - y_2' = C\bar{B}(\bar{A} - I)(u_1 - u_2).$$
	This vanishes for all $u_1 \neq u_2$ only if $\bar{A} = I$, which degenerates the SSM to a memoryless map and eliminates all temporal memory - contradicting the purpose of state-space modelling. In the selective Mamba SSM, the input-dependent operators $\bar{A}_k(s_k)$ and $\bar{B}_k(s_k)$ introduce additional asymmetric cross-terms that further break permutation symmetry. $\square$
	
	Bidirectionality alone does not resolve the problem. In the backward pass, $\text{Rev}_t(h^{\leftarrow})$ reverses along the \emph{temporal} axis $L'$, not along the modality-sequence axis $\sum C_i'$. The backward SSM still processes modality blocks in the order $M, M{-}1, \dots, 1$; each modality's backward hidden state therefore accumulates context from those that follow it in the reversed sequence. The averaged representation
	$$H = \frac{h^{\rightarrow} + \text{Rev}_t(h^{\leftarrow})}{2}$$
	remains order-dependent: both the forward and backward context seen by modality $i$ change qualitatively with its position in the arbitrary concatenation, so the same physical modality receives a structurally different representation depending solely on where it was placed in $m$.
	
	\subsection{Differentiable Modality Ordering via Gumbel-Sinkhorn}
	
	We replace the fixed concatenation order with a \emph{learned, input-dependent permutation} computed differentiably via the Gumbel-Sinkhorn operator \cite{mena2018}. The key intuition is that the optimal ordering is itself a function of the data: for a sample in which audio is uninformative, placing it first in the forward SSM pass deposits unreliable context into the hidden states of all subsequent modalities, whereas placing it last confines the contamination to the final position.
	
	\paragraph{Uniform channel projection.}
	Different InitEncoders naturally produce output representations of different widths: for instance, an audio encoder may output $C_{\text{audio}}' = 64$ channels while a video encoder outputs $C_{\text{video}}' = 256$ channels. These \emph{heterogeneous dimensions} prevent treating the modality outputs as interchangeable objects of the same type, which is a prerequisite for learning a permutation over them. To resolve this, a modality-specific affine map $f_i : \R^{C_i'} \to \R^{C'}$, defined by $f_i(v) = W_i^{\text{proj}} v + b_i^{\text{proj}}$ with $W_i^{\text{proj}} \in \R^{C' \times C_i'}$ and $b_i^{\text{proj}} \in \R^{C'}$, is appended to each InitEncoder to linearly project its output into a shared embedding space of dimension $C'$. After this projection, $u_i \in \R^{L' \times C'}$ for all $i \in \{1,\dots,M\}$. The matrix $U$ is then the concatenation of these uniformly-dimensioned representations along the channel axis: $U = u_1 \circ u_2 \circ \dots \circ u_M \in \R^{L' \times MC'}$, where each $u_i$ contributes exactly $C'$ columns.
	
	\paragraph{Score matrix.}
	A \emph{meta-network} $\phi$ is a small auxiliary network whose sole purpose is to predict how modalities should be ordered, operating on top of the main feature representations without participating in the primary task of engagement prediction. Concretely, $\phi$ maps the set $\{u_i\}_{i=1}^M$ to an $M \times M$ assignment \emph{score} matrix $Z \in \R^{M \times M}$, where the entry $Z_{ij}$ quantifies the fitness of placing modality $j$ at sequence position $i$ in the inter-modal BiMamba: a large value of $Z_{ij}$ means the meta-network believes modality $j$ should be processed at step $i$. Each modality representation is first summarised by temporal mean-pooling:
	$$\bar{u}_i = \frac{1}{L'}\sum_{t=1}^{L'} u_i^{(t)} \in \R^{C'}.$$
	Two asymmetric projections into a score-embedding space of dimension $d_s$ then produce query and key vectors:
	$$q_i = W_Q\,\bar{u}_i + b_Q \in \R^{d_s}, \qquad k_j = W_K\,\bar{u}_j + b_K \in \R^{d_s},$$
	where $W_Q, W_K \in \R^{d_s \times C'}$ and $b_Q, b_K \in \R^{d_s}$ are learnable parameters shared across modalities. The score is then computed as a scaled dot product:
	$$Z_{ij} = \frac{q_i^\top k_j}{\sqrt{d_s}}.$$
	The asymmetry of the query-key factorisation is intentional: it allows the model to distinguish between a modality acting as a \emph{context supplier} (the key role) and a modality acting as a \emph{context receiver} (the query role), and assigns independent scores to each direction. A symmetric formulation $Z_{ij} = z_i^\top z_j$ would force $Z = Z^\top$, which in turn constrains the learned permutation matrix $P$ to be symmetric. The only symmetric permutation matrices are \emph{involutions} - permutations that are their own inverse ($\pi = \pi^{-1}$), corresponding to products of disjoint transpositions (swaps). This would unnecessarily halve the space of learnable orderings, excluding, for example, cyclic permutations of three or more modalities.
	
	\paragraph{Gumbel perturbation.}
	The Gumbel distribution is classically used in extreme value theory to model the asymptotic distribution of the maximum of a sequence of independent and identically distributed random variables. Its probability density function for location $\mu$ and scale $\beta$ is given by: 
	$$f(x) = \frac{1}{\beta} e^{-(x-\mu)/\beta} e^{-e^{-(x-\mu)/\beta}}$$ 
	In deep learning, the Gumbel-Max trick \cite{jang2017} leverages this distribution to draw discrete samples from a categorical distribution in a fully differentiable manner. By adding independent Gumbel noise to unnormalized log-probabilities (scores), the deterministic argmax of these perturbed values becomes mathematically equivalent to drawing a sample from the true underlying multinomial distribution. 
	
	To generate standard Gumbel noise ($\mu=0, \beta=1$) programmatically, we apply inverse transform sampling. This technique maps a standard uniform random variable to a target distribution by passing it through the inverse of that target's cumulative distribution function (CDF). Given that the standard Gumbel CDF is $F(x) = e^{-e^{-x}}$, setting a uniform random variable $U \sim \mathrm{Uniform}(0,1)$ equal to $F(x)$ and solving explicitly for $x$ yields the inverse mapping $x = -\ln(-\ln U)$. Thus, to enable stochastic exploration of orderings during training, the score matrix $Z$ is perturbed with i.i.d.\ standard Gumbel noise and scaled by a temperature parameter $\tau > 0$:
	$$\tilde{Z}_{ij} = \frac{Z_{ij} + G_{ij}}{\tau}, \qquad G_{ij} = -\ln(-\ln U_{ij}), \quad U_{ij} \overset{\text{i.i.d.}}{\sim} \mathrm{Uniform}(0,1).$$
	The temperature $\tau$ is annealed over the course of training to progressively concentrate probability mass towards a single hard permutation. We adopt a standard exponential decay schedule, updating the temperature at each epoch $t$ via $\tau_t = \max(\tau_\infty, \tau_0 \cdot \gamma^t)$, where $\gamma \in (0, 1)$ represents the decay rate, starting from an initial exploration temperature $\tau_0$ (e.g., $1.0$) down to a tight categorical approximation threshold $\tau_\infty$ (e.g., $0.05$).
	
	\paragraph{Sinkhorn normalisation.}
	The \textbf{Birkhoff-von Neumann theorem} \cite{birkhoff1946} establishes that the set of doubly-stochastic matrices (square matrices with non-negative real entries whose rows and columns each sum to exactly $1$) constitutes a convex polytope whose extreme points are precisely the set of discrete permutation matrices. Geometrically, a convex polytope is a bounded, multi-dimensional solid defined by flat facets, and its extreme points represent its absolute vertices. Because any interior point of a convex polytope can be expressed uniquely as a convex combination (a weighted average) of its vertices, any doubly-stochastic matrix can be decomposed into a convex combination of hard permutation matrices. This structural property makes the Birkhoff polytope a mathematically principled continuous relaxation of discrete permutations.
	
	A standard Softmax operator is insufficient in this context because it operates independently across a single dimension, ensuring only that the rows sum to $1$ (yielding a row-stochastic matrix). In an inter-modal fusion framework where rows represent sequence positions and columns map to distinct modalities, a row-stochastic matrix introduces severe structural redundancies. It permits multiple sequence positions to concurrently assign high probability mass to the exact same dominant modality column (such as audio), which causes the network to attend to duplicate representations while completely suppressing and ignoring other available modalities. 
	
	Constraining both rows and columns to sum to $1$ eliminates this issue by enforcing a strict continuous one-to-one assignment map. Following \cite{mena2018}, we obtain a valid doubly stochastic matrix $P$ by applying the iterative Sinkhorn-Knopp algorithm directly to the elementwise exponential of the perturbed score matrix:
	$$P = \text{Sinkhorn}\bigl(\exp(\tilde{Z})\bigr).$$
	The Sinkhorn-Knopp algorithm achieves normalisation by alternating between row and column scaling operations. Because the exponential map $\exp(\tilde{Z})$ guarantees that all initial entries are strictly positive ($>0$), the resulting row and column sums are bounded strictly away from zero. This structural guarantee completely eliminates the risk of numerical division-by-zero errors (NaN), ensuring that the iterative scaling steps remain entirely smooth and differentiable. Consequently, the gradient operator $\displaystyle \frac{\partial P}{\partial Z}$ remains mathematically well-conditioned throughout optimization, enabling robust, end-to-end backpropagation of the latent modal ordering.
	
	\paragraph{Soft reordering.}
	With $P$ computed as above, the permuted representation for position $i$ is the $P$-weighted convex combination of all intra-modal outputs. Concretely, the $i$-th row of $P$, denoted $P_{i\cdot} = (P_{i1}, \dots, P_{iM})$, defines a probability distribution over the $M$ modalities; since $P$ is doubly-stochastic, we have $P_{ij} \geq 0$ and $\sum_j P_{ij} = 1$ for every $i$. The soft assignment is then:
	$$\tilde{u}_i = \sum_{j=1}^{M} P_{ij}\,u_j \in \R^{L' \times C'}, \qquad i \in \{1,\dots,M\}.$$
	The entry $P_{ij} \in [0,1]$ is the soft probability that modality $j$ is assigned to position $i$ in the inter-modal sequence; when $P$ is close to a hard permutation matrix (as happens after temperature annealing), exactly one $P_{ij}$ per row will be close to 1 and the rest close to 0, recovering a near-discrete assignment. The fixed concatenation $m$ is replaced by
	$$\tilde{m} = \text{Transpose}(\tilde{u}_1 \circ \tilde{u}_2 \circ \dots \circ \tilde{u}_M) \in \R^{MC' \times L'},$$
	and the inter-modal BiMamba operates on $\tilde{m}$ instead of $m$; all equations in Section 1.2.2 remain unchanged under this substitution.
	
	\paragraph{Inference.}
	At test time, Gumbel noise is suppressed ($G_{ij} = 0$). If a hard permutation is required (e.g.\ for latency-critical deployment), the optimal discrete assignment is recovered by solving the linear assignment problem over the score matrix $Z$:
	$$\pi^* = \arg\min_{\pi \in S_M} \sum_{i=1}^{M} \bigl(-Z_{i,\pi(i)}\bigr),$$
	where $S_M$ denotes the symmetric group - the set of all $M!$ bijections from $\{1,\dots,M\}$ to itself, i.e.\ all possible orderings of the $M$ modalities. This minimisation is solved by applying the Hungarian algorithm to $-Z$ in $O(M^3)$ time, which is negligible for the small modality counts typical in multimodal sensing.
	
	\paragraph{Equivariance guarantee.}
	Let $\rho \in S_M$ be applied to the inputs before the meta-network. Because $Z_{ij}$ is computed from the data via the shared query-key projections, relabelling the inputs by $\rho$ relabels the rows and columns of $Z$ by $\rho$ as well, and consequently the recovered hard permutation transforms as $\pi^*_\rho = \rho \circ \pi^*$. The permuted sequence fed to the inter-modal BiMamba is therefore the same physical ordering of modalities regardless of how the practitioner originally numbered them. Consequently, the fused representation $H = \text{BiMamba}(\tilde{m})$ is invariant to the practitioner's arbitrary choice of concatenation order, provided the model has converged to a stable optimal ordering $\pi^*$.
	
	\newpage
	\begin{thebibliography}{9}
		
		\bibitem{ziyu2025}
		Ziyu Jia et al., \emph{A Multimodal BiMamba Network with Test-Time Adaptation for Emotion Recognition Based on Physiological Signals}, 39th Conference on Neural Information Processing Systems (NeurIPS 2025).
		
		\bibitem{xie2026mhc}
		Zhenda Xie et al., \emph{mHC: Manifold-Constrained Hyper-Connections}, arXiv preprint arXiv:2512.24880v2 (2026).
		
		\bibitem{gu2023mamba}
		Albert Gu and Tri Dao, \emph{Mamba: Linear-Time Sequence Modeling with Selective State Spaces}, arXiv preprint arXiv:2312.00752 (2024).
		
		\bibitem{mena2018}
		Gonzalo Mena et al., \emph{Learning Latent Permutations with Gumbel-Sinkhorn Networks}, International Conference on Learning Representations (ICLR 2018).
		
		\bibitem{jang2017}
		Eric Jang, Shixiang Gu, and Ben Poole, \emph{Categorical Reparameterization with Gumbel-Softmax}, International Conference on Learning Representations (ICLR 2017).
		
		\bibitem{gumbel1954}
		Emil Julius Gumbel, \emph{Statistical Theory of Extreme Values and Some Practical Applications}, National Bureau of Standards Applied Mathematics Series, 33 (1954).
		
		\bibitem{birkhoff1946}
		Garrett Birkhoff, \emph{Tres observaciones sobre el algebra lineal}, Universidad Nacional de Tucum\'{a}n, Revista, Serie A, 5:147--151 (1946).
		
	\end{thebibliography}
	
\end{document}